\chapter{Introduction}
\label{ch:introduction}

\section{Context and Motivation}
\label{sec:introduction-context-motivation}

% Introduce eBPF in generale: origine, evoluzione, ruolo nel kernel Linux.
The extended Berkeley Packet Filter (eBPF) is a new programmability mechanism that changes the way we can observe and interact with the Linux kernel. Originally it was designed as a packet filtering mechanism, but now the paradigm has evolved to a more general purpose, addressing the need for observability and security in the Linux kernel.

% Spiega il trade-off: estendere il kernel senza modificarlo, ma con rischi.
eBPF introduces a tradeoff in kernel development: on one hand, it allows developers to extend the kernel with new features without the need to modify the kernel source code; but on the other hand, this widely adopted Linux kernel technology presents a new challenge for security, because it allows developers to run custom code in kernel space, which is a privileged environment.

% Introduci verifier e sicurezza senza entrare troppo nei dettagli.
When we use this mechanism to load custom programs in the kernel, we need to be sure that the code is safe and does not compromise the security of the system. 
eBPF addresses this problem by introducing a sophisticated verifier framework. This framework is able to analyze the code before it is loaded in the kernel, and to reject any code that does not respect the security and performance constraints.
It permits programs running in the user space to submit eBPF bytecode to the kernel, which is then verified. 

% Spiega perche' eBPF e' adatto a monitoraggio runtime.
Once the programs are verified, they are executed in the kernel space, achieving high performance and low overhead through Just-in-Time (JIT) compilation, while maintaining or even enhancing the security of the kernel.

% Collega eBPF agli hook kernel.
eBPF allows attaching eBPF programs to specific kernel hooks, enabling the observation and collection of events in the kernel space, without the need to modify the kernel source code.

% Introduci la pipeline della tesi.
This capability enables a new approach to security monitoring, allowing us to observe security-relevant events in a reliable and structured way, collecting them from different kernel hooks, transporting them to the user space, decoding them and making them usable for further analysis.
As a result, the impact of this new technology is significant, because it allows us to build a complete pipeline from the kernel space to the user space.

% Chiudi con impatto e motivazione generale.
Organizations leverage this new technology, because it permits to leverage the combination of safety, performance and observability, to build next-generation security tools that run in a privileged environment.

\section{Thesis Objectives}
\label{sec:introduction-thesis-objectives}

% Presenta l'obiettivo principale della tesi.
This thesis investigates the use of eBPF for security monitoring through the design and implementation of a prototype system that leverages eBPF to collect relevant events from the Linux kernel, transports them to user space, and converts them into structured information for further analysis.% Metti una bullet list con gli obiettivi specifici.

In order to achieve this goal, the thesis will focus on the following specific objectives: \begin{itemize}
    % [TODO: aggiungi obiettivi specifici, RICORDATI DI AGGIUNGERE LE POLICIES]
    \item Design and implement a kernel-side component that uses eBPF to attach to selected kernel hooks and collect raw relevant events from the kernel space.
    \item Implement a user-space component that receives the raw events from kernel space, decodes their context and arguments, enriches selected fields, and finally outputs them in a structured format that can be used for further analysis.
    \item Define an event model and decoding pipeline that separates kernel-side data collection from user-space interpretation and presentation.
    \item Evaluate the prototype in terms of functional correctness, event coverage, and runtime overhead on the target Linux environment.
\end{itemize}

% Chiarisci cosa e' implementato e cosa resta futuro.
% [TODO: scope attuale vs future work]

\section{Scope of the Work}
\label{sec:introduction-scope}

% Chiarisci che il lavoro e' un prototipo di tesi, non un prodotto completo.
The work presented in this thesis is a research prototype that investigates the use of eBPF for runtime security monitoring. Its scope is to demonstrate how kernel-level events can be collected, transported to user space, decoded, and exposed in a structured form.

% Spiega il rapporto con Tracee.
\textit{Tracee} is used as the main architectural reference, since it represents a production-grade example of eBPF-based runtime security monitoring. Unlike production-grade tools that must support a wide range of kernels and deployment scenarios, this prototype is intentionally tailored to a specific target environment, namely \textit{Rocky Linux 8.10}.

% Delimita il focus tecnico.
The focus of this work is the implementation of kernel hooks related to process and security events, in particular a selected set of events relevant to process execution, filesystem activity, credential changes, memory operations, namespaces, and kernel-hardening signals.

% Chiarisci cosa resta fuori o in evoluzione.
The resulting architecture is intended to provide a basis for further extensions, including policy-based filtering, anomaly detection, and event correlation.

\section{Thesis Structure}
\label{sec:introduction-thesis-structure}

% Scrivi questa sezione alla fine.
The remainder of this thesis is organized as follows: \begin{itemize}
    \item \textbf{Chapter 2} introduces the technical background required to understand the work and discusses related runtime security tools, with particular attention to eBPF architecture, the Linux kernel, the Verifier, the eBPF hooks, and the Tracee runtime security tool.
    \item \textbf{Chapter 3} presents the requirements and the design of the proposed system, including the target environment, the main architectural components, and the design choices that shaped the prototype.
    \item \textbf{Chapter 4} describes the implementation of the system, covering the eBPF kernel-side component, the Go user-space runtime, event transport, decoding, filtering, output generation, and the main implementation challenges.
    \item \textbf{Chapter 5} evaluates the prototype in terms of functional behavior, event coverage, runtime overhead, and limitations observed on the target environment.
    \item \textbf{Chapter 6} concludes the thesis by summarizing the achieved results, discussing the current limitations, and outlining future extensions.
\end{itemize}

